%% kapa: exam=1 round=35 year=2024 period=2 form=A subject=appraisal_law subject_ko=감정평가관계법규 seq=16 number=16
%% source: Q-Net 감정평가사 기출문제 (cst00309) artlSeq=5213229 file=raw/1008/5213229/2교시 문제지 원본.pdf
%% conversion: kordoc 4.24.51 -> tools/gen1.py (auto); review=proofread-sonnet; flags=none
\begin{question}{16}
부동산 가격공시에 관한 법령상 지가의 공시에 관한 설명으로 옳은 것은?
\begin{choices}
  \item 개별공시지가에 이의가 있는 자는 그 결정ㆍ공시일부터 60일 이내에 서면으로 관할 관청에 이의를 신청할 수 있다.
  \item 표준지공시지가의 단위면적은 3.3제곱미터로 한다.
  \item 개발부담금의 부과대상이 아닌 토지에 대하여는 개별공시지가를 결정ㆍ공시하여야 한다.
  \item 표준지공시지가의 공시에는 표준지에 대한 지목 및 용도지역이 포함되어야 한다.
  \item 개별공시지가의 결정ㆍ공시에 드는 비용은 30퍼센트 이내에서 국고에서 보조한다.
\end{choices}
\end{question}
\answer{4}
